\documentclass[12pt]{article}
\usepackage[utf8]{inputenc}
\usepackage[T1]{fontenc}
\usepackage[margin=1in]{geometry}
\usepackage{setspace}
\usepackage{xcolor}
\usepackage[colorlinks=true, urlcolor=blue!60!black, linkcolor=blue!60!black]{hyperref}
\setlength{\parindent}{2em}
\setlength{\parskip}{1.25em}
\renewcommand{\baselinestretch}{1.0}
\begin{document}
\begin{center}
{\LARGE \textit{CFPB v. Community Financial Services Ass'n of America, Ltd.}, 601 U.S. 416 (2024)\par}
\vspace{0.8em}
{\large Case Brief --- Four-Page Edition\par}
\end{center}
\vspace{1em}

\section{Memory Jogger}

Rejecting payday lenders' challenge, the Court held 7--2 that the CFPB's standing funding mechanism---drawing a capped amount each year from the Federal Reserve System's combined earnings---satisfies the Appropriations Clause, because an appropriation need only be a law identifying a source of public funds and authorizing their expenditure for designated purposes (\href{https://supreme.justia.com/cases/federal/us/601/22-448/}{Justia}).

\section{Detailed Case Facts}

After the 2008 financial crisis, the Dodd-Frank Act created the Consumer Financial Protection Bureau within the Federal Reserve System, consolidating 18 consumer-finance statutes and giving it rulemaking, enforcement, and adjudicatory power over ``unfair, deceptive, or abusive'' practices. Instead of annual appropriations, the Director may requisition each year from ``the combined earnings of the Federal Reserve System'' the amount the Director determines ``to be reasonably necessary to carry out'' the Bureau's duties, subject to an inflation-adjusted cap of 12\% of the System's 2009 operating expenses (about \$734 million in FY 2022). The Bureau may retain and invest unspent funds. \textit{Seila Law} (2020) had already severed the Director's removal protection. In 2017 the CFPB issued a Payday Lending Rule restricting repeated withdrawal attempts from borrowers' accounts.

\section{Procedural History}

Two payday-lending trade associations challenged the rule in the Western District of Texas, arguing among other things that the funding scheme violated the Appropriations Clause. The district court granted summary judgment to the CFPB. The Fifth Circuit reversed and vacated the rule, holding the Bureau's ``self-actualizing, perpetual funding mechanism'' unconstitutional. The Court granted certiorari, heard argument on October 3, 2023, and decided the case on May 16, 2024.

\section{Judicial Votes}

Reversed and remanded, \textbf{7--2}.
\begin{itemize}
\item \textbf{Majority:} Justice Thomas, joined by Chief Justice Roberts and Justices Sotomayor, Kagan, Kavanaugh, Barrett, and Jackson.
\item \textbf{Concurrences:} Justice Kagan, joined by Justices Sotomayor, Kavanaugh, and Barrett; Justice Jackson.
\item \textbf{Dissent:} Justice Alito, joined by Justice Gorsuch.
\end{itemize}

\section{Holding}

The CFPB's funding mechanism satisfies the Appropriations Clause. An ``Appropriatio[n] made by Law'' requires only a statute that identifies a specified source of public funds and authorizes their expenditure for designated purposes. The Clause does not require annual renewal, line-item specificity, or a fixed sum rather than a cap.

\section{Analysis of Opinions}

\textbf{Majority (Thomas).} Because the Federal Reserve's surplus would otherwise go to the Treasury, the funding is ``Money \ldots\ drawn from the Treasury'' subject to the Clause. Founding-era dictionaries defined an appropriation as setting apart public funds for designated uses, and the Army Clause's two-year limit shows by negative implication that other appropriations need not be time-limited. Parliament enacted standing grants and ``any Sum not exceeding'' caps, and the First Congress enacted capped lump sums and standing, fee-based funding for the Customs Service and Post Office. The Bureau does not set its own budget: Congress fixed the cap, and discretion to draw \textit{less} mirrors early practice. The Clause, phrased as a limitation, imposes no freestanding demand for annual review.

\textbf{Kagan (concurring).} Two centuries of ``long settled and established practice''---lump sums, standing appropriations funding nearly two-thirds of the FY 2022 budget, and assessment-funded bank regulators---show the CFPB's mechanism ``would have fit right in'' at every point.

\textbf{Jackson (concurring).} Where the text imposes no limit on Congress's institutional design, courts should not invent one to second-guess Congress's response to the financial crisis.

\textbf{Alito (dissenting).} ``Appropriations'' is a term of art forged in parliamentary resistance to the Stuarts. The CFPB's combination of perpetuity, self-selected draws, indirect funding through the Federal Reserve, retained surpluses, exemption from appropriations review, and vast power is unprecedented, and \textit{Seila Law} made it worse by placing this fiscally autonomous agency under presidential control.

\section{Examples: Future Applications}

\begin{enumerate}
\item \textbf{Same-Side (Valid): Fee-Funded Registry.} A permanent statute letting a trademark registry spend its filing fees up to a cap identifies a source and purpose, like the early Customs Service and Post Office.
\item \textbf{Same-Side (Valid): Permanent ``Such Sums'' Benefit.} A statute paying veterans a stipend from the Treasury indefinitely names a source and purpose; Kagan noted that such standing appropriations fund programs like Social Security and Medicare.
\item \textbf{Opposite-Side (Invalid): Standing Army Draw.} A permanent, capped draw for an Army office violates the Army Clause's two-year limit, on which the majority relied.
\item \textbf{Opposite-Side (Invalid): Spending Without a Statute.} Paying contractors \$200 million for a program Congress refused to fund, on Article II alone, lacks any appropriation, as in \textit{OPM v. Richmond}.
\item \textbf{Fence-Sitter: Uncapped AI Agency.} An agency drawing ``whatever amounts the Director deems necessary'' from the Treasury in perpetuity names a source and purpose, but challengers can stress the majority's reliance on the CFPB's cap and raise the nondelegation question the Court left open.
\end{enumerate}

\section{Critique}

\textbf{A toothless test?} Critics echoing Justice Alito argue that the ``source plus purpose'' test is formalistic and nearly toothless, apparently allowing a perpetual, uncapped appropriation for any lawful executive purpose and gutting Congress's check of the purse.

\textbf{A defense.} Defenders answer that any future Congress can repeal or amend a standing appropriation by ordinary legislation, and that Alito's multi-factor ``novelty'' test lacks a textual anchor and would require arbitrary line-drawing about how much fiscal independence is too much.

\section{Key Quotations}

\begin{quote}
``Under the Appropriations Clause, an appropriation is simply a law that authorizes expenditures from a specified source of public money for designated purposes.''
\end{quote}
\noindent --- Justice Thomas, for the Court, 601 U.S. at 424.

\begin{quote}
``Based on the Constitution's text, the history against which that text was enacted, and congressional practice immediately following ratification, we conclude that appropriations need only identify a source of public funds and authorize the expenditure of those funds for designated purposes to satisfy the Appropriations Clause.''
\end{quote}
\noindent --- Justice Thomas, for the Court, 601 U.S. at 426.

\begin{quote}
``Unfortunately, today's decision turns the Appropriations Clause into a minor vestige. The Court upholds a novel statutory scheme under which the powerful Consumer Financial Protection Bureau (CFPB) may bankroll its own agenda without any congressional control or oversight. \ldots\ It is not an exaggeration to say that the CFPB enjoys a degree of financial autonomy that a Stuart king would envy.''
\end{quote}
\noindent --- Justice Alito, dissenting, 601 U.S. at 448, 467.

\end{document}
